% TEMPLATE-MILC-v3.tex - plantilla maestra UNICA de las guias MILC v3.
%
% La usan las cuatro familias de guias, cada una con su driver:
%   · Grados 6-11          -> scripts/build-guias-g11.py
%   · Bebras               -> scripts/build-guias-bebras.py
%   · Semillero            -> scripts/build-guias-semillero.py
%   · Territorio interior  -> scripts/build-guias-territorio-interior.py
%
% Antes habia tres copias casi identicas (~400 lineas iguales, 6 distintas):
% cada mejora habia que aplicarla tres veces, y en la practica no se hacia
% (el motor de recursos vivia solo en dos de ellas). Lo que varia entre
% familias esta parametrizado en 6 placeholders que cada driver llena
% (se nombran aqui SIN su sintaxis real: el builder reemplaza por texto plano,
% tambien dentro de los comentarios, y un valor multilinea rompe el preambulo):
%
%   HEADER_IZQ        encabezado izquierdo de pagina
%   HEADER_DER        encabezado derecho de pagina
%   PORTADA_ETIQUETA  rotulo sobre el numero grande (GRADO/MOMENTO/MODULO)
%   PORTADA_SERIE     linea de serie de la portada
%   PORTADA_ID        identificador de la guia en la portada
%   PORTADA_CIRCULO   el circulito de grado (°), o vacio si no aplica
%   PDF_TITULO        titulo en texto plano (sin \textbf ni comillas TeX) para
%                     los metadatos del PDF (/Title); lo produce lib_xelatex.texto_pdf
%
% Composicion (auditoria 2026-09-03): las bandas de estacion piden espacio
% (\Needspace*) y prohiben el salto justo despues (\nopagebreak), asi no quedan
% huerfanas al pie; las cajas largas (softbox, drawbox, infoband, puentebox) son
% `breakable`, asi una caja de media pagina ya no arrastra todo a la siguiente
% dejando la anterior al 40 %. Todo texto sobre mostaza o verde va en milcNegro
% (blanco sobre #E5B400 da 1,93:1 y sobre #58A923 2,95:1; ninguno pasa WCAG AA).
% El resto de claves las documenta content/guias/_SCHEMA.md. El script valida
% que no queden placeholders sin reemplazar antes de escribir el archivo final.

% Etiquetado PDF (latex-lab): /Lang, estructura y texto alternativo de las
% figuras, para lectores de pantalla. Debe ir ANTES de \documentclass.
\DocumentMetadata{lang=es-CO,tagging=on}
\documentclass[11pt,letterpaper]{article}
% headheight=24pt: la cabecera derecha lleva dos lineas (identidad de la guia
% y estacion actual). headsep se acorta en la misma medida para que el bloque
% de texto quede exactamente donde estaba con headheight=12pt.
\usepackage[margin=2.0cm,headheight=24pt,headsep=13pt]{geometry}
\usepackage[table]{xcolor}
\usepackage{fontspec}
\usepackage[spanish]{babel}
\usepackage{tikz}
% Librerías TikZ para los diagramas de las guías (bloque `recursos:`):
%  · babel  → babel en español vuelve activos caracteres como «>», y TikZ los usa
%             en sus opciones (>=stealth, ->). Sin esta librería, cualquier
%             diagrama con flechas revienta con «\language@active@arg> has an
%             extra }». Debe ir ANTES de que se dibuje nada.
%  · positioning        → sintaxis «below left=2pt and 3pt».
%  · decorations.pathreplacing → llaves ({brace}) para señalar rangos.
\usetikzlibrary{babel,positioning,decorations.pathreplacing}
\usepackage{graphicx}
% Circuitos: el trabajo de electrónica se hace hoy con micro:bit y sus sensores
% integrados (MakeCode, sin cableado), así que no se necesitan esquemáticos.
% Si algún día entran montajes con protoboard, basta descomentar esta línea:
% los diagramas .tex/.tikz declarados en `recursos:` ya se incrustan solos.
% \usepackage{circuitikz}
\usepackage[most]{tcolorbox}
\usepackage{tabularx}
\usepackage{array}
\usepackage{fancyhdr}
\usepackage{titlesec}
\usepackage[document]{ragged2e}  % bandera a la derecha en todo el cuerpo
\usepackage{enumitem}
\usepackage{microtype}
\usepackage{etoolbox}
\usepackage{needspace}
% hyperref al final del preambulo: metadatos del PDF (titulo, autor, idioma,
% familia), marcadores por estacion (\pdfbookmark) y enlaces sin recuadro.
\usepackage[hidelinks,
  pdftitle={Sostenerse afuera: la red no aparece, se arma},
  pdfauthor={PhD. Álvaro Cárdenas Orozco},
  pdfsubject={Territorio interior · Educación socioemocional · Grado 11° · Momento 6 · MILC v3},
  pdflang={es-CO},
  bookmarksopen=true,bookmarksnumbered=false]{hyperref}

% Helvetica Neue no trae la flecha U+2192, asi que cada «→» escrito literalmente
% en el YAML se compilaba a NADA, en silencio (462 flechas en 61 guias: rutas del
% tipo «paleta → tipografia → lockup» salian rotas). Mapeamos el caracter a su
% version matematica, que si tiene glifo. Asi el autor puede seguir escribiendo
% «→» comodamente en el YAML.
\usepackage{newunicodechar}
\newunicodechar{→}{\ensuremath{\rightarrow}}
% Mismo problema con los simbolos de marca y vinetas: el «✓» de «una palomita ✓»
% salia como cuadrito vacio en el PDF de todas las guias que lo usaban.
\usepackage{pifont}
\newunicodechar{✓}{\ding{51}}
\newunicodechar{✗}{\ding{55}}
\newunicodechar{✔}{\ding{52}}
\newunicodechar{•}{\textbullet}
\newunicodechar{≥}{\ensuremath{\geq}}
\newunicodechar{≤}{\ensuremath{\leq}}

\setmainfont{Helvetica Neue}
\setsansfont{Helvetica Neue}
\setlength{\parindent}{0pt}
\setlength{\parskip}{6pt}
% Sin particion de palabras: en 6.º-7.º una palabra cortada por guion al final
% de linea («Comuni-cacion») frena la lectura mucho mas que una linea corta.
\hyphenpenalty=10000
\exhyphenpenalty=10000
\pretolerance=2000
\tolerance=3000
\setlength{\emergencystretch}{3em}
\linespread{1.10}
\renewcommand{\arraystretch}{1.25}
\setlist{leftmargin=5mm,itemsep=1mm,topsep=1mm}
\newcolumntype{Y}{>{\RaggedRight\arraybackslash}X}

\definecolor{milcMagenta}{HTML}{E6007E}
\definecolor{milcVino}{HTML}{5A0038}
\definecolor{milcTurquesa}{HTML}{0093A5}
\definecolor{milcVerde}{HTML}{58A923}
\definecolor{milcMostaza}{HTML}{E5B400}
\definecolor{milcOcre}{HTML}{B86600}
\definecolor{milcNegro}{HTML}{241816}
\definecolor{milcGris}{HTML}{F7F7F4}
\definecolor{milcLinea}{HTML}{E8E2DD}
\definecolor{milcGradePrimary}{HTML}{0D9488}
\definecolor{milcGradeDark}{HTML}{115E59}
\definecolor{milcGradeSoft}{HTML}{CCFBF1}
\definecolor{milcGradeAccent}{HTML}{CFE7FF}
\definecolor{milcGradeText}{HTML}{FFFFFF}
\color{milcNegro}

\pagestyle{fancy}
\fancyhf{}
% Cabecera de dos lineas: identidad de la guia arriba y, a la derecha y en
% gris, la estacion en curso (\markright desde \milcestacion). Antes de la
% primera estacion la segunda linea va vacia; el \strut mantiene la altura
% para que la linea de arriba no baile entre paginas.
\lhead{\small Territorio interior · Educación socioemocional\\[-1pt]{\footnotesize\strut}}
\rhead{\small Grado 11° · Momento 6 · MILC v3\\[-1pt]{\footnotesize\strut\color{milcNegro!65}\rightmark}}
\cfoot{\small\thepage}
\renewcommand{\headrulewidth}{0pt}

% Sobre mostaza (#E5B400) y verde (#58A923) el texto va en milcNegro; sobre el
% resto de la paleta, en blanco. Se usa dentro de fonttitle/fontupper, donde un
% \color posterior manda sobre coltitle.
\newcommand{\milctintasobre}[1]{%
  \ifboolexpr{test{\ifstrequal{#1}{milcMostaza}} or test{\ifstrequal{#1}{milcVerde}}}%
    {\color{milcNegro}}{\color{white}}}

\newtcolorbox{titlebox}[1]{
  enhanced,colback=#1,colframe=#1,arc=7mm,boxrule=0pt,
  left=7mm,right=7mm,top=3mm,bottom=3mm,
  before skip=8pt,after skip=8pt,
  fontupper=\sffamily\bfseries\Large\color{white}
}

\newtcolorbox{softbox}[3]{
  enhanced,breakable,pad at break=3mm,
  colback=#2,colframe=#1,borderline west={4pt}{0pt}{#1},
  arc=2mm,boxrule=.4pt,left=4mm,right=4mm,top=3mm,bottom=3mm,
  before skip=6pt,after skip=8pt,title={#3},coltitle=white,
  colbacktitle=#1,fonttitle=\sffamily\bfseries\milctintasobre{#1},fontupper=\RaggedRight,
  attach boxed title to top left={xshift=3mm,yshift=-2mm},
  boxed title style={arc=2mm, boxrule=0pt}
}

\newtcolorbox{drawbox}[2]{
  enhanced,breakable,pad at break=4mm,
  colback=white,colframe=#1,arc=4mm,boxrule=1pt,
  left=5mm,right=5mm,top=4mm,bottom=4mm,before skip=8pt,after skip=8pt,
  title={#2},coltitle=white,colbacktitle=#1,fonttitle=\sffamily\bfseries\milctintasobre{#1},
  fontupper=\RaggedRight,
  attach boxed title to top left={xshift=4mm,yshift=-2mm},
  boxed title style={arc=3mm, boxrule=0pt}
}

\newtcolorbox{infoband}[1]{
  enhanced,breakable,pad at break=2mm,colback=#1!8,colframe=#1!50,arc=2mm,boxrule=.5pt,
  left=4mm,right=4mm,top=2mm,bottom=2mm,before skip=4pt,after skip=4pt,
  fontupper=\small\RaggedRight
}

% Puente narrativo entre secciones (contrato editorial Conéctate).
% Da continuidad: "Acabas de X. Ahora vas a Y porque Z."
% Visualmente: banda izquierda en color del grado, texto en itálica pequeña.
\newtcolorbox{puentebox}{
  enhanced,breakable,pad at break=2mm,colback=milcGradeSoft,colframe=milcGradeDark,
  borderline west={3pt}{0pt}{milcGradeDark},
  arc=0mm,boxrule=0pt,left=5mm,right=4mm,top=2.5mm,bottom=2.5mm,
  before skip=4pt,after skip=8pt,
  fontupper=\itshape\small\color{milcNegro}\RaggedRight
}
% La flecha va en modo matematico a proposito: Helvetica Neue NO tiene el glifo
% U+2192, asi que \textrightarrow se compilaba a NADA («Missing character: There
% is no → in font Helvetica Neue») y el puente salia sin flecha. En modo
% matematico la flecha viene de la fuente matematica y si se dibuja.
\newcommand{\puente}[1]{%
  \begin{puentebox}%
    \textcolor{milcGradeDark}{$\rightarrow$}\hspace{2mm}#1%
  \end{puentebox}%
}

\newcommand{\linead}[1]{\noindent\color{milcLinea}\rule{#1}{0.5pt}\color{milcNegro}}
\newcommand{\respuesta}[1]{\par\vspace{2mm}\foreach \n in {1,...,#1}{\noindent\color{milcLinea}\rule{\linewidth}{0.5pt}\par\vspace{5mm}}\color{milcNegro}}
\newcommand{\checkbox}{\textcolor{milcGradeDark}{\fbox{\phantom{x}}}\hspace{5pt}}

% ─── Listas aireadas para las guías socioemocionales ────────────────────
% Componer las verificaciones como «\checkbox texto\\» deja el texto que salta
% de línea debajo del cuadrito y apelmaza el bloque. Estos entornos alinean la
% continuación y airean los ítems. Los usa Territorio Interior; cualquier
% familia puede hacerlo.
\newlist{checklista}{itemize}{1}
\setlist[checklista]{
  label=\textcolor{milcGradeDark}{\fbox{\phantom{x}}},
  leftmargin=9mm, labelsep=3.2mm, itemsep=2.4mm,
  topsep=2.5mm, parsep=0pt, partopsep=0pt, before=\RaggedRight
}
\newlist{pasoslista}{enumerate}{1}
\setlist[pasoslista]{
  label=\textcolor{milcGradeDark}{\sffamily\bfseries\arabic*.},
  leftmargin=8mm, labelsep=3mm, itemsep=2.8mm,
  topsep=1mm, parsep=0pt, partopsep=0pt, before=\RaggedRight
}

\newcommand{\phasecard}[4]{
\begin{tcolorbox}[enhanced,colback=#3,colframe=#2,arc=3mm,boxrule=.5pt,height=3.05cm,valign=center,halign=center,left=1.5mm,right=1.5mm,top=2mm,bottom=2mm]
{\sffamily\bfseries\fontsize{22}{24}\selectfont\textcolor{#2}{#1}}\par
{\sffamily\bfseries\small #4}
\end{tcolorbox}
}

% ── Piezas del rediseno 2026 (legibilidad 6.º-11.º) ───────────────────
% Banda de estacion: reemplaza el titlebox macizo. Titulo a la izquierda,
% dato practico (tiempo/modalidad) a la derecha, en una sola linea.
% Nunca huerfana: pide 9 lineas libres (o salta de pagina rellenando con
% \vfill) y prohibe el salto justo despues de la banda. Nueve y no seis:
% la caja partible que sigue necesita ~80pt (titulo adosado, relleno y dos
% lineas) para arrancar en la misma pagina; con menos, tcolorbox la manda
% entera a la siguiente y la banda se queda sola. El penalty 10000 va
% en `after`, ANTES del espacio posterior: un `after skip` (aunque sea 0pt)
% es un \vskip tras la caja, y ese glue es un punto de corte legal que un
% \nopagebreak escrito despues de \end{tcolorbox} ya no cierra. Cada banda
% deja marcador PDF y actualiza la cabecera derecha.
\newcounter{milcestacion}
\newcommand{\milcestacion}[2]{%
\Needspace*{9\baselineskip}%
\stepcounter{milcestacion}%
\pdfbookmark[1]{#2}{milcestacion\themilcestacion}%
\markright{#2}%
\begin{tcolorbox}[enhanced,colback=#1,colframe=#1,arc=2mm,boxrule=0pt,
  left=5mm,right=5mm,top=2.6mm,bottom=2.6mm,before skip=11pt,
  after={\par\nopagebreak\vskip7pt}]
{\sffamily\bfseries\fontsize{15}{18}\selectfont\milctintasobre{#1}#2}
\end{tcolorbox}}

% Tarjeta de paso numerada: sustituye las tablas «Pilar / Aplicacion», que
% eran formato de consulta docente, no de estudiante. El numero va en un
% circulo sobre el borde para que la secuencia se lea de un vistazo.
\newcommand{\milcpaso}[3]{%
\Needspace{4\baselineskip}%
\begin{tcolorbox}[enhanced,colback=white,colframe=milcLinea,arc=1.5mm,boxrule=.9pt,
  left=14mm,right=4mm,top=3mm,bottom=3mm,before skip=4pt,after skip=4pt,
  fontupper=\RaggedRight,
  overlay={\node[anchor=north west,circle,fill=#2,text=white,inner sep=0pt,
    minimum size=7.4mm,font=\sffamily\bfseries\small]
    at ([xshift=3.6mm,yshift=-3.2mm]frame.north west) {#1};}]
#3
\end{tcolorbox}}

% Casilla marcable de tamano util con lapiz (la anterior era \fbox{\phantom{x}},
% de alto variable segun el texto que la seguia).
\newcommand{\milcmarca}{\framebox[3.8mm]{\rule{0pt}{3.4mm}}}

% Fila marcable de lista suelta (checks de la estacion 1).
\newcommand{\milccheckrow}[1]{%
\par\vspace{1.8mm}\noindent
\makebox[6.4mm][l]{\raisebox{-.55mm}{\textcolor{milcNegro}{\milcmarca}}}%
\parbox[t]{\dimexpr\linewidth-6.4mm\relax}{\RaggedRight #1}\par}

% Celda marcable dentro de la rubrica (tabla de autoevaluacion). Misma
% sangria francesa, pero pensada para vivir dentro de una columna X.
\newcommand{\milcopcion}[1]{%
\makebox[5.2mm][l]{\raisebox{-.55mm}{\textcolor{milcNegro}{\milcmarca}}}%
\parbox[t]{\dimexpr\linewidth-5.2mm\relax}{\RaggedRight #1}}

% ── Recursos visuales de la guía ──────────────────────────────────────
% Declarados en el bloque `recursos:` del YAML y emitidos por el builder.
% \guiaFigura   → imagen rasterizada o PDF (foto, carta celeste, montaje).
% \guiaDiagrama → fuente TikZ/circuitikz (.tex) incrustada como vector:
%                 es la vía para circuitos y esquemáticos editables.
% [#1] = texto alternativo (alt) · #2 = ruta relativa al .tex · #3 = pie
% El alt llega al PDF etiquetado (/Alt) y lo lee el lector de pantalla.
\newcommand{\guiaFigura}[3][]{%
\begin{center}
\includegraphics[width=0.88\linewidth,height=0.38\textheight,keepaspectratio,alt={#1}]{#2}\\[1.5mm]
{\footnotesize\itshape\textcolor{milcNegro!75}{#3}}
\end{center}
}

\newcommand{\guiaDiagrama}[2]{%
\begin{center}
\input{#1}\\[1.5mm]
{\footnotesize\itshape\textcolor{milcNegro!75}{#2}}
\end{center}
}

\begin{document}
\thispagestyle{empty}

% ── Portada ───────────────────────────────────────────────────────────
% Antes: un rectangulo de color a pagina completa. Se veia bien en pantalla,
% pero en la fotocopiadora del colegio se comia el toner y salia gris sucio.
% Ahora la banda ocupa el tercio superior y el resto va en blanco.
\begin{tikzpicture}[remember picture,overlay]
  \path[fill=milcGradePrimary]
    (current page.north west) rectangle ([yshift=-10.6cm]current page.north east);

  \node[anchor=north west,text=milcGradeText] at ([xshift=2.0cm,yshift=-1.55cm]current page.north west)
    {\sffamily\bfseries\fontsize{12}{14}\selectfont MOMENTO};
  \node[anchor=north west,text=milcGradeText,opacity=.85] at ([xshift=2.0cm,yshift=-2.35cm]current page.north west)
    {\sffamily\bfseries\fontsize{8.5}{10}\selectfont TERRITORIO INTERIOR · SOCIOEMOCIONAL};
  \node[anchor=north east,text=milcGradeText,opacity=.85,align=right] at ([xshift=-2.0cm,yshift=-1.55cm]current page.north east)
    {\sffamily\bfseries\fontsize{9}{12}\selectfont I.E. Sor María Juliana\\Cartago, Valle del Cauca};

  \node[anchor=west,text=milcGradeText] (gradonum) at ([xshift=1.85cm,yshift=-7.1cm]current page.north west)
    {\sffamily\bfseries\fontsize{112}{112}\selectfont 6};
  % El circulito de grado (°) solo aplica a las guias de grado («11°»); en
  % Bebras y Semillero el numero grande es un momento/modulo, no un grado.
  % Se posiciona relativo al nodo `gradonum`, no en coordenadas absolutas.
  

  \node[anchor=west,text=milcGradeText,text width=11.4cm,align=left] at ([xshift=1.5cm,yshift=0.5cm]gradonum.east)
    {
     {\sffamily\bfseries\fontsize{9}{11}\selectfont\MakeUppercase{Territorio interior · Socioemocional}}\\[3mm]
     {\sffamily\bfseries\fontsize{21}{25}\selectfont\setlength{\baselineskip}{27pt}Sostenerse\\[4pt]afuera\\[4pt]la red no aparece, se arma}\\[3.5mm]
     {\sffamily\bfseries\fontsize{15}{18}\selectfont Grado 11° · Momento 6}
    };
\end{tikzpicture}

\vspace*{9.4cm}
\vfill

{\sffamily\bfseries\fontsize{8}{10}\selectfont\textcolor{milcGradeDark}{LO QUE TE LLEVAS HOY}}

\begin{tcolorbox}[enhanced,colback=white,colframe=milcNegro,arc=2mm,boxrule=1.6pt,
  left=5mm,right=5mm,top=4mm,bottom=4mm,before skip=3pt,after skip=12pt,
  fontupper=\RaggedRight\fontsize{11}{15.5}\selectfont]
Tu plan de red: el inventario de los cuatro lugares donde hoy te encuentras con gente sin proponértelo, la lista de los que van a desaparecer al graduarte, tres actividades concretas con horario fijo para el año que viene, y las dos personas a las que te comprometes a escribirles.
\end{tcolorbox}

\begin{tcolorbox}[enhanced,colback=milcGris,colframe=milcGris,arc=2mm,boxrule=0pt,
  left=5mm,right=5mm,top=3.5mm,bottom=3.5mm,after skip=12pt]
{\sffamily\bfseries\fontsize{8}{10}\selectfont\textcolor{milcNegro!55}{ESTA GUÍA ES DE}}\\[3mm]
\begin{tabularx}{\linewidth}{X p{3.2cm} p{3.2cm}}
\linead{\linewidth} & \linead{3.0cm} & \linead{3.0cm}\\[-1mm]
{\fontsize{8.5}{10}\selectfont\textcolor{milcNegro!55}{Nombre completo}} &
{\fontsize{8.5}{10}\selectfont\textcolor{milcNegro!55}{Grupo}} &
{\fontsize{8.5}{10}\selectfont\textcolor{milcNegro!55}{Fecha}}\\
\end{tabularx}
\end{tcolorbox}

\vfill

\noindent\textcolor{milcLinea}{\rule{\linewidth}{1pt}}\\[2mm]
{\sffamily\bfseries\fontsize{9.5}{12}\selectfont Autor: Dr. Álvaro Cárdenas Orozco}\\[1mm]
{\sffamily\fontsize{8.5}{11}\selectfont\textcolor{milcNegro!55}{Metodología MILC · Apertura ancestral · Soledad y adultez temprana · Triángulo Dussel-Estoicismo-Floridi}}

\newpage

\section*{Apertura: Sostenerse afuera: la red no aparece, se arma}

\begin{softbox}{milcGradeDark}{milcGradeSoft}{Saber ancestral}
En \textbf{1986}, bajo el rectorado de \textbf{Harold Rizo Otero}, la Universidad del Valle creó su \textbf{sistema de regionalización}: llevar la universidad a los municipios en vez de exigirle a todo el mundo que se fuera a Cali. \textbf{Las primeras sedes fueron Buenaventura, Buga, Caicedonia, Palmira, Tuluá, Zarzal, Sevilla y \emph{Cartago}}, y el \textbf{20 de octubre de 1986} empezaron los programas presenciales. \emph{La sede de Cartago venía de antes:} \textbf{había arrancado en modalidad a distancia con apoyo del SENA, que aportó la infraestructura entre 1984 y 1985.} \emph{Y hay una manera de leer esto que sirve para esta guía.} \textbf{Alguien miró un problema ---la gente tenía que \emph{irse} para estudiar--- y en vez de pedirle a cada joven que fuera más valiente, \textbf{cambió la estructura}:} \emph{puso la universidad donde estaba la gente.} \textbf{Fíjate en la lógica, porque es la de todo este momento:} \emph{cuando algo le pasa a mucha gente a la vez ---irse solo, quedarse sin red, no encontrar dónde estudiar---, casi nunca es un problema de carácter}. \textbf{Es un problema de estructura, y las estructuras se pueden construir.} \emph{Cartago se llama «ciudad bisagra» por su posición ---entre el Valle y el eje cafetero---, y una bisagra es exactamente eso: lo que permite que dos cosas se conecten.}
\end{softbox}

\begin{softbox}{milcTurquesa}{milcGris}{Saber contemporáneo}
¿Y qué tanto importa tener gente alrededor? \textbf{Más de lo que cualquiera diría, y está medido con una precisión incómoda.} \textbf{Julianne Holt-Lunstad}, \textbf{Timothy Smith} y \textbf{Bradley Layton} reunieron \textbf{148 estudios} con datos de \textbf{308.849 personas}, seguidas en promedio \textbf{7,5 años}. \textbf{El resultado:} \emph{las personas con relaciones sociales adecuadas tienen un \textbf{50\,\% más de probabilidad de supervivencia}} que quienes tienen relaciones pobres o insuficientes (Holt-Lunstad, Smith y Layton, 2010). \textbf{Y la comparación que hicieron los autores es la que hay que retener:} \emph{la magnitud de ese efecto es \textbf{comparable a dejar de fumar}, y \textbf{supera} a factores de riesgo tan conocidos como la obesidad y el sedentarismo.} \emph{El resultado se mantuvo teniendo en cuenta la edad, el sexo, el estado de salud inicial, el tiempo de seguimiento y la causa de muerte.} \textbf{Traducido:} \emph{tener con quién contar no es un lujo emocional ni un tema de «ser sociable» ---es un factor de salud del mismo orden que los que sí nos enseñan a cuidar---.} \textbf{Y a nadie le enseñan a construirlo.}
\end{softbox}

\begin{softbox}{milcMostaza}{milcGris}{Pregunta-puente}
\textit{Alguien resolvió un problema de mucha gente \textbf{cambiando la estructura, no pidiendo más valentía}. \textbf{Cuando te gradúes, ¿qué estructura vas a perder} \ldots \textbf{y con qué la vas a reemplazar}?}
\end{softbox}

\begin{softbox}{milcVerde}{milcGris}{Saber y saber hacer}
Al terminar podrás: (1) \textbf{reconocer} que el colegio venía haciendo por ti el trabajo de ponerte con gente todos los días, y que eso se acaba; (2) \textbf{entender} que la red social pesa en la salud tanto como los factores de riesgo clásicos; (3) \textbf{identificar} qué hace que un lugar produzca vínculos ---repetición, tiempo compartido, algo que hacer juntos---; (4) \textbf{planear} tres actividades con horario fijo para el año que viene y comprometerte a sostener dos vínculos concretos.
\end{softbox}



\small
\begin{tabularx}{\textwidth}{>{\bfseries}p{3.0cm}Y}
\rowcolor{milcGradePrimary!12}Autor & Dr. Álvaro Cárdenas Orozco\\
DBA & Comprende los fenómenos que afectan a su territorio, regula sus emociones ante ellos y acompaña a otros con empatía (transversal 6°--11°, articulado con la política «Escuela, territorio de vida» del MEN, Resolución 006519 de 2025, y la Ley 1620 de 2013)\\
Referentes & Primera ayuda psicológica (OMS, 2011/2012) · Directrices IASC (2007) · USGS y Servicio Geológico Colombiano · Pueblo nasa y la minga (CRIC) · Dussel · Estoicismo · Floridi\\
Producto final & Tu plan de red: el inventario de los cuatro lugares donde hoy te encuentras con gente sin proponértelo, la lista de los que van a desaparecer al graduarte, tres actividades concretas con horario fijo para el año que viene, y las dos personas a las que te comprometes a escribirles.\\
\end{tabularx}
\normalsize

\puente{En el momento 5 miraste si irte o quedarte. \textbf{Hoy da igual cuál escojas:} \emph{en las dos, la estructura que te rodeaba doce años se acaba en noviembre}.}

\milcestacion{milcGradeDark}{Ruta MILC: así vamos a aprender}
\begin{tabularx}{\textwidth}{>{\centering\arraybackslash}X>{\centering\arraybackslash}X>{\centering\arraybackslash}X>{\centering\arraybackslash}X}
\phasecard{1}{milcVerde}{milcGris}{Escucha\\[-1mm]\small Observo, escucho y pregunto.} &
\phasecard{2}{milcTurquesa}{milcGris}{Sistematización\\[-1mm]\small Armo lo que antes venía dado.} &
\phasecard{3}{milcMagenta}{milcGris}{Praxis\\[-1mm]\small Construyo y aplico.} &
\phasecard{4}{milcVino}{milcGris}{Evaluación\\[-1mm]\small Reflexión liberadora.}
\end{tabularx}


\puente{Primero, un inventario sencillo y revelador: dónde te encuentras con gente sin haberlo planeado.}

\milcestacion{milcVerde}{Estación 1 de 3 · Escucha}

\begin{softbox}{milcVerde}{milcGris}{Escena para escuchar}
\textbf{Actividad 1 · IDENTIFICA --- Dónde me encuentro con gente.} \textbf{Nadie va a leer esta página.} \textbf{Primera parte.} Haz la lista de \textbf{los lugares donde te encuentras con gente sin proponértelo}: \emph{el salón, el descanso, la ruta, el entrenamiento, la iglesia, el barrio, un grupo de música, el trabajo si trabajas.} \textbf{Segunda parte, la que revela.} Al frente de cada uno marca: \emph{¿va a seguir existiendo cuando te gradúes?} \textbf{Sé literal ---la mayoría no---.} \textbf{Tercera parte.} Escribe \textbf{los nombres de las cinco personas con las que más hablas} y responde, para cada una: \emph{¿dónde la conociste? ¿la seguirías viendo si no fuera por ese lugar?} \textbf{Y la última:} \emph{¿cuándo fue la última vez que hiciste un amigo nuevo, y dónde?} \emph{Si te cuesta acordarte, ese es el dato: los amigos que tienes se hicieron casi todos en lugares que alguien más organizó.}
\end{softbox}

{\sffamily\bfseries\fontsize{8}{10}\selectfont\textcolor{milcNegro!55}{MARCA CUANDO LO HAYAS HECHO}}
\milccheckrow{Listaste los lugares donde te encuentras con gente sin proponértelo.}
\milccheckrow{Marcaste cuáles \textbf{van a seguir existiendo} después de graduarte.}
\milccheckrow{Anotaste dónde conociste a las cinco personas con las que más hablas.}

\begin{infoband}{milcVerde}


\textbf{En tu cuaderno (Actividad 1 · IDENTIFICA).} Título: «Actividad 1. Dónde me encuentro con gente». \textbf{Formato:} los lugares con la marca de si sobreviven + las cinco personas y dónde las conociste + cuándo hiciste un amigo nuevo por última vez. \textbf{Extensión:} media página. \textbf{Sabes que terminaste cuando} sabes \textbf{cuántos de tus lugares desaparecen}. Esta página es tuya: \textbf{nadie más tiene que leerla si tú no quieres}.
\end{infoband}

\puente{Ya lo tienes. Ahora, qué venía haciendo el colegio por ti, cuánto pesa la red y qué hace que un lugar produzca vínculos.}

\milcestacion{milcTurquesa}{Fase 2 · Sistematización: entender la red}

\begin{softbox}{milcTurquesa}{milcGris}{Cuatro claves sobre sostenerse}
Cuatro claves para entender por qué a tanta gente le pega el primer año fuera del colegio.
\end{softbox}

\milcpaso{1}{milcTurquesa}{\textbf{El colegio venía haciendo un trabajo enorme y gratis.} \emph{Durante doce años, alguien te puso todos los días, a la misma hora, en el mismo salón, con las mismas cuarenta personas, haciendo cosas juntos.} \textbf{Esa es, literalmente, la receta para producir vínculos ---y no la hiciste tú---.} \emph{Por eso la gente cree que «los amigos se hacen solos»:} \textbf{se hacían solos porque la estructura estaba puesta.} \textbf{Y en noviembre se acaba de golpe.} \emph{No gradualmente ---de un día para otro---.} \textbf{Ese es el motivo de que a tanta gente le pegue el primer año afuera y crea que le pasa algo raro:} \emph{no le pasa nada raro. Perdió una infraestructura y nadie le avisó que existía.}}
\milcpaso{2}{milcTurquesa}{\textbf{La red pesa en la salud como los factores que sí nos enseñan a cuidar.} \emph{148 estudios, 308.849 personas, 7,5 años de seguimiento:} \textbf{quienes tienen relaciones sociales adecuadas tienen un 50\,\% más de probabilidad de supervivencia}, \emph{y el efecto es \textbf{comparable a dejar de fumar} y \textbf{mayor} que el de la obesidad y el sedentarismo} (Holt-Lunstad y otros, 2010). \textbf{Detente en la comparación.} \emph{A todo el mundo le enseñan que fumar mata y que hay que moverse.} \textbf{A nadie le enseñan que quedarse sin gente hace un daño del mismo orden}, \emph{y por eso nadie planea su red como planea su carrera.} \emph{Esto no es un argumento sentimental:} \textbf{es una razón de salud para tratar la construcción de vínculos como algo que se hace a propósito.}}
\milcpaso{3}{milcTurquesa}{\textbf{Los vínculos necesitan repetición, no química.} \emph{Mira tu Actividad 1:} \textbf{casi todos tus amigos salieron de lugares donde te tocaba volver}. \emph{No de un encuentro afortunado ---de haber vuelto el martes siguiente---.} \textbf{Un lugar produce vínculos cuando tiene tres cosas:} \emph{la misma gente} ---no gente distinta cada vez---, \emph{con frecuencia} ---semanal, no mensual--- \emph{y \textbf{algo que hacer juntos}}, que da de qué hablar sin tener que inventar conversación. \textbf{Por eso una clase produce amistades y una fiesta casi nunca}, \emph{y por eso un equipo, un coro, un grupo de estudio o un voluntariado funcionan mejor que «salir a conocer gente».} \emph{La consecuencia práctica es antipática y muy útil:} \textbf{no se buscan amigos ---se buscan actividades con horario fijo, y los amigos salen de ahí---.}}
\milcpaso{4}{milcTurquesa}{\textbf{Pedir es una habilidad, y a muchos les enseñaron lo contrario.} \emph{Acuérdate del momento 2 de este grado y del mandato del grado 9:} \textbf{a los hombres sobre todo, pero no solo a ellos, se les enseñó que pedir ---ayuda, compañía, un favor--- es de débiles.} \emph{Y en la adultez temprana esa lección se cobra:} \textbf{la gente se queda sola sabiendo que hay personas a quienes podría escribirle, y no escribe.} \emph{Las razones son siempre las mismas:} \textbf{«no quiero molestar», «ya no tenemos nada en común», «hace mucho no hablamos», «si le importara me escribiría él».} \emph{Y hay un dato que conviene tener:} \textbf{casi todo el mundo se alegra cuando alguien le escribe después de mucho tiempo}, \emph{y casi todo el mundo cree que va a ser el único al que le da pena hacerlo.}}

\begin{drawbox}{milcTurquesa}{Qué hace que un lugar produzca vínculos}
\begin{pasoslista}
\item \textbf{La misma gente} vuelve. \emph{No sirve un lugar con gente distinta cada vez.}
\item \textbf{Con frecuencia:} semanal funciona; mensual casi no.
\item \textbf{Con algo que hacer juntos:} da conversación sin tener que inventarla.
\item \textbf{Durante meses:} las amistades no se hacen en tres semanas.
\item \textbf{Y la regla que resume todo:} \emph{no busques amigos, busca un horario.}
\end{pasoslista}
\end{drawbox}

\begin{softbox}{milcMostaza}{milcGris}{Errores comunes y su corrección}
\textbf{«Los amigos se hacen solos»:} se hacían solos porque el colegio ponía la estructura. \textbf{Creer que el bajón del primer año es personal:} le pasa a casi todo el mundo y tiene causa estructural. \textbf{Buscar amigos en vez de actividades:} lo que produce vínculos es el horario. \textbf{Fiestas y eventos únicos:} sin repetición no producen nada. \textbf{«No quiero molestar»:} casi todo el mundo se alegra de recibir un mensaje. \textbf{Esperar a que el otro escriba:} los dos están esperando. \textbf{Confundir estar en línea con estar acompañado:} no es lo mismo y el cuerpo lo nota.
\end{softbox}

\begin{infoband}{milcTurquesa}


\textbf{En tu cuaderno (Actividad 2 · EXPLICA).} Título: «Actividad 2. Lo que hacía el colegio». \textbf{Formato:} explica qué estructura te daba el colegio y por qué su desaparición no es un problema de carácter; después escribe las \textbf{cuatro condiciones} que hacen que un lugar produzca vínculos. \textbf{Extensión:} media página. \textbf{Sabes que terminaste cuando} puedes explicar la comparación de \textbf{Holt-Lunstad con dejar de fumar}.
\end{infoband}

\puente{Ya sabes qué se pierde. Es hora de reemplazarlo con algo concreto: horarios, no intenciones.}

\milcestacion{milcMagenta}{Fase 3 · Praxis: armar la red}

\begin{softbox}{milcMagenta}{milcGris}{Cómo se arma una red a propósito}
Esto no se resuelve con buenas intenciones ni con ser más sociable. Se resuelve con tres horarios y dos mensajes.
\end{softbox}

\milcpaso{1}{milcMagenta}{\textbf{Los cuatro lugares (6 min).} Termina el inventario y \textbf{pásalo por las cuatro condiciones} ---misma gente, frecuencia, algo que hacer, meses---. \emph{Vas a ver cuáles de tus lugares actuales cumplen las cuatro,} \textbf{y esos son el modelo de lo que hay que reemplazar.}}
\milcpaso{2}{milcMagenta}{\textbf{Lo que desaparece (6 min).} Escribe \textbf{la lista de lo que se acaba en noviembre} y, al frente de cada uno, \emph{a quién dejarías de ver si no hicieras nada}. \textbf{Con nombres.} \emph{Esta lista incomoda y por eso funciona:} \textbf{casi todo el mundo descubre que entre cinco y quince personas dependen enteramente de una estructura que va a desaparecer.}}
\milcpaso{3}{milcMagenta}{\textbf{Tres actividades con horario (12 min).} Escribe \textbf{tres actividades} para el año que viene que cumplan \textbf{las cuatro condiciones}. \emph{Ejemplos reales:} \textbf{un equipo o un entrenamiento, un grupo de estudio fijo, un coro o una banda, un voluntariado, la Junta de Acción Comunal ---donde puedes afiliarte desde los catorce, como verás en el momento 7---, un semillero, un curso, un trabajo con compañeros estables.} \textbf{Con día y hora tentativos.} \emph{Y una condición dura:} \textbf{al menos una tiene que ser presencial.}}
\milcpaso{4}{milcMagenta}{\textbf{Las dos personas (8 min, y esta semana).} Escoge \textbf{dos personas} con las que quieras seguir después de noviembre ---no las cinco más cercanas: piensa en quién se va a perder si nadie hace nada--- y escribe \textbf{cómo vas a sostener cada vínculo}: \emph{cada cuánto, por qué medio, y quién escribe primero ---que vas a ser tú---.} \textbf{Y hazlo concreto:} \emph{«el primer domingo de cada mes le escribo» funciona; «mantener el contacto» no.} \emph{Después escríbele hoy a una de las dos.}}

\begin{drawbox}{milcMagenta}{Antes de cerrar tu plan de red}
\begin{checklista}
\item Pasaste tus lugares actuales por las \textbf{cuatro condiciones}.
\item Tienes la lista de \textbf{lo que desaparece} con nombres de quién se perdería.
\item Tus tres actividades cumplen las cuatro condiciones y tienen \textbf{día y hora}.
\item Al menos una es \textbf{presencial}.
\item Tus dos vínculos tienen \textbf{frecuencia, medio y quién escribe primero}.
\item \textbf{Ya le escribiste} a una de las dos personas.
\end{checklista}
\end{drawbox}

\begin{softbox}{milcMagenta}{milcGris}{Plantilla de trabajo}
\textbf{Las cuatro condiciones.} \emph{«¿Misma gente? · ¿Semanal? · ¿Algo que hacer juntos? · ¿Durante meses?»} \\[1mm]
\textbf{Mi actividad.} \emph{«\_\_\_\_, los \_\_\_\_ a las \_\_\_\_, desde \_\_\_\_.»} \\[1mm]
\textbf{El vínculo.} \emph{«A \_\_\_\_ le escribo cada \_\_\_\_. Escribo yo primero.»}
\end{softbox}

\begin{infoband}{milcMagenta}


\textbf{En tu cuaderno (Actividad 3 · CREA).} Título: «Actividad 3. Mi plan de red». \textbf{Formato:} el inventario con las cuatro condiciones + la lista de lo que desaparece con nombres + las tres actividades con día y hora + los dos vínculos con su frecuencia + a quién le escribiste. \textbf{Extensión:} una página. \textbf{Sabes que terminaste cuando} \textbf{ya mandaste un mensaje}.
\end{infoband}

\puente{El producto tiene horarios adentro. \emph{Los amigos no se hacen con ganas: se hacen con repetición.}}

\milcestacion{milcMostaza}{Producto de la sesión}
\textbf{Plan de red: lo que desaparece al graduarse, tres actividades con horario fijo y dos vínculos sostenidos a propósito}.
\begin{softbox}{milcMostaza}{milcGris}{Criterios de calidad}
(1) El inventario evalúa los lugares actuales contra las cuatro condiciones que producen vínculos. (2) La lista de lo que desaparece nombra a las personas que se perderían sin acción. (3) Las tres actividades cumplen las cuatro condiciones, tienen día y hora, y al menos una es presencial. (4) Los dos vínculos tienen frecuencia, medio y responsable de iniciar el contacto. (5) Se envió efectivamente un primer mensaje, no solo se planeó.
\end{softbox}

\puente{Antes de cerrar, revisa lo que hace fracasar estos planes: si tus actividades no tienen día y hora fijos, no van a producir a nadie.}

\milcestacion{milcVino}{Evaluación liberadora · cinco dimensiones}

\begin{softbox}{milcVino}{milcGris}{Sentido de la evaluación}
Evaluar no es solo revisar una nota. Es reconocer qué aprendí, qué cambió en mí, qué emociones manejé, cómo me relaciono con la ciudad y la comunidad, y qué le dejaré a quien viene detrás.
\end{softbox}

\textbf{Mi reflexión liberadora en cinco dimensiones.} Completa cada frase iniciada:

\small
\begin{tabularx}{\textwidth}{>{\bfseries\RaggedRight\arraybackslash}p{3.4cm}Y>{\RaggedRight\arraybackslash}p{4.6cm}}
\rowcolor{milcVino}\color{white}Dimensión & \color{white}\bfseries Mi respuesta (frame starter) & \color{white}\bfseries Algo que descubrí\\
1. Personal & Lo que más cambió en mí fue \linead{6.5cm} & \linead{4.2cm}\\
2. Emocional & La emoción más fuerte fue \linead{2.5cm} y la regulé \linead{3cm} & \linead{4.2cm}\\
3. Ciudadana & En democracia esto importa porque \linead{6cm} & \linead{4.2cm}\\
4. Local (Cartago) & En mi barrio sería útil para \linead{6.5cm} & \linead{4.2cm}\\
5. Intergeneracional & A mi abuela/abuelo le diría \linead{6.5cm} & \linead{4.2cm}\\
\end{tabularx}
\normalsize

\begin{infoband}{milcVino}
\textbf{Para la próxima sustentación.} Vuelve a esta tabla en seis meses. Verás que las respuestas cambian — no porque lo aprendido cambie, sino porque tú cambias. La evaluación liberadora es una conversación contigo a través del tiempo.
\end{infoband}


\puente{Cerramos con tres voces: la comunidad que se funda escuchando, la abeja y el enjambre, y la compañía que hoy se confunde con conexión.}

\milcestacion{milcGradeDark}{Triángulo de pensamiento · cierre filosófico}

\begin{softbox}{milcVino}{milcGris}{Dussel · lente del nosotros}
\textit{«Aquel que es capaz de escuchar silenciosamente al Otro como otro es el que puede constituir una comunidad y no una mera sociedad totalizada.»}

{\footnotesize\itshape\color{milcNegro!70}--- Enrique Dussel · Filosofía de la liberación (1977)}

\emph{Fíjate en el verbo: \textbf{constituir}.} \textbf{Una comunidad no se encuentra ---se constituye, es decir, alguien la hace---.} \emph{Y eso es exactamente lo contrario de lo que uno espera a los diecisiete, cuando la comunidad venía puesta y parecía natural.} \textbf{La regionalización de Univalle es un ejemplo del mismo verbo a otra escala:} \emph{alguien decidió que la universidad fuera adonde estaba la gente, y ochenta municipios cambiaron por eso.} \textbf{Nadie encontró esa oportunidad: alguien la constituyó.} \emph{Y con tu red va a pasar igual ---no vas a \emph{encontrar} un grupo: vas a tener que volver a un lugar varias semanas seguidas hasta que exista---.}

\textbf{¿Estoy esperando encontrar gente, o estoy dispuesto a volver a un mismo sitio doce semanas seguidas?}
\end{softbox}
\linead{\linewidth}\\[1mm]
\linead{\linewidth}

\begin{softbox}{milcMostaza}{milcGris}{Estoicismo · Marco Aurelio · Meditaciones VI, 54 (c. 175 d.C.) · cuidado interior}
\textit{«Lo que no beneficia al enjambre, tampoco beneficia a la abeja.»}

{\footnotesize\itshape\color{milcNegro!70}--- Marco Aurelio · Meditaciones VI, 54 (c. 175 d.C.)}

\emph{Esta frase se lee siempre en una dirección ---lo que le sirve al grupo---.} \textbf{Aquí interesa la otra:} \emph{una abeja sola no es una abeja más libre. Es una abeja que no dura.} \textbf{Y eso, que suena a metáfora, resultó ser un dato:} \emph{las relaciones sociales adecuadas se asocian con un 50\,\% más de probabilidad de supervivencia, con un efecto comparable a dejar de fumar} (Holt-Lunstad y otros, 2010). \emph{A los diecisiete, la independencia suele imaginarse como no necesitar a nadie.} \textbf{La independencia real es otra cosa:} \emph{poder escoger de quién dependes, y saber construir esa dependencia a propósito.} \textbf{Nadie se sostiene solo, y el que lo intenta no es más fuerte ---está en riesgo---.}

\textbf{Cuando pienso en ser independiente, ¿me imagino sin necesitar a nadie?}
\end{softbox}
\linead{\linewidth}\\[1mm]
\linead{\linewidth}

\begin{softbox}{milcTurquesa}{milcGris}{Floridi · lente de la infoesfera}
\textit{«El florecimiento de las entidades informacionales y de toda la infoesfera debe promoverse preservando, cultivando y enriqueciendo sus propiedades.»}

{\footnotesize\itshape\color{milcNegro!70}--- Luciano Floridi · La ética de la información (2013; trad.)}

\textbf{Hay una confusión que esta generación paga cara: creer que estar conectado es estar acompañado.} \emph{Se puede pasar el día entero hablando con gente y llegar la noche sin haber visto a nadie,} \textbf{y el cuerpo nota la diferencia aunque la cabeza no la registre.} \emph{No es que lo digital no sirva ---sirve mucho para \textbf{sostener} lo que ya existe, y por eso el plan de hoy incluye dos personas y una frecuencia---.} \textbf{Lo que casi no hace es \emph{crear} vínculos nuevos}, \emph{porque le faltan las cuatro condiciones: la misma gente, la frecuencia obligada, algo que hacer juntos y meses.} \textbf{Regla práctica para el año que viene:} \emph{lo digital para mantener a los de antes; lo presencial para que aparezcan los nuevos.}

\textbf{Ayer hablé con mucha gente. ¿A cuántas vi?}
\end{softbox}
\linead{\linewidth}\\[1mm]
\linead{\linewidth}

\begin{infoband}{milcNegro}
\textbf{Nota para el docente.} Las tres son citas textuales y llevan su referencia debajo. Puedes remitir a la obra en clase.
\end{infoband}

\milcestacion{milcVino}{Mi autoevaluación · marca una casilla por fila}

{\small
\setlength{\extrarowheight}{2.2pt}
\begin{tabularx}{\textwidth}{>{\RaggedRight\arraybackslash\bfseries}p{3.0cm}YYY}
\rowcolor{milcVino}\color{white}\bfseries Criterio & \color{white}\bfseries Lo logré & \color{white}\bfseries En proceso & \color{white}\bfseries Necesito apoyo\\[.6mm]
Comprensión del tema & \milcopcion{Explico las ideas con mis palabras.} & \milcopcion{Reconozco las ideas, falta precisión.} & \milcopcion{Necesito repasar lo básico.}\\[.8mm]
\rowcolor{milcGris}Calidad del plan & \milcopcion{Sé qué estructura desaparece al graduarme y tengo tres actividades con horario fijo para reemplazarla.} & \milcopcion{Reconozco que voy a perder gente, pero creo que los amigos se hacen solos.} & \milcopcion{Sigo pensando que pedir ayuda o compañía es de débiles.}\\[.8mm]
Aplicación y producto & \milcopcion{Mi producto cumple los criterios.} & \milcopcion{Está armado, con ajustes pendientes.} & \milcopcion{Aún está en construcción.}\\[.8mm]
\rowcolor{milcGris}Reflexión 5 dimensiones & \milcopcion{Respondí cada una con honestidad.} & \milcopcion{Respondí algunas con detalle.} & \milcopcion{Mis respuestas son muy generales.}\\[.8mm]
Triángulo de pensamiento & \milcopcion{Conecté las 3 voces con mi trabajo.} & \milcopcion{Conecté al menos una.} & \milcopcion{No las relaciono con mi proceso.}\\[.8mm]
\end{tabularx}}

\vspace{3mm}
\textbf{Hoy entendí que:}\\[1mm]
\linead{\linewidth}\\[1mm]
\linead{\linewidth}

\textbf{Mi compromiso para la próxima sesión es:}\\[1mm]
\linead{\linewidth}\\[1mm]
\linead{\linewidth}

\begin{softbox}{milcMostaza}{milcGris}{Cita de despedida · Marco Aurelio}
\textit{``El obstáculo en el camino se vuelve el camino. Lo que estorba al camino se vuelve camino.''}\\[1mm]
Cada error de hoy, bien observado, se convierte en pieza del camino que pisarás mañana.
\end{softbox}

\small
\begin{tabularx}{\textwidth}{>{\bfseries}p{3.6cm}Y}
\rowcolor{milcGradeDark!12}Nombre completo: & \linead{10cm}\\
Fecha: & \linead{4cm} \quad Curso/grupo: \linead{4cm}\\
Mi firma: & \linead{9cm}\\
\end{tabularx}
\normalsize

\begin{infoband}{milcGradeDark}
\textbf{Para la próxima sesión.} Dos cosas. \textbf{Primero, escríbele a una de tus dos personas} y trae \textbf{qué respondió}. \emph{Y averigua de verdad una de tus tres actividades: dónde queda, a qué hora es, si hay que inscribirse, cuánto cuesta. No basta con nombrarla ---la mitad de estos planes se caen porque nadie averiguó el horario---.} \textbf{Segundo, pregunta en tu casa por el primer año afuera:} averigua con alguien mayor \textbf{cómo fue su primer año después del colegio} ---si se fue o se quedó, con quién contaba, cuánto tardó en hacer amigos nuevos, si la pasó sola--- y \textbf{qué le habría servido saber}. \textbf{Trae:} la respuesta a tu mensaje y lo que te contaron. \emph{Prepárate, porque esta conversación suele sorprender: mucha gente que hoy se ve rodeada cuenta que ese primer año fue de los más solos de su vida, y casi nadie lo dijo en su momento. Saber que es común no lo evita, pero evita lo peor, que es creer que a uno le pasa algo raro.}
\end{infoband}

\end{document}
